\pagetopsection{ Creating a Witch}


\begin{multicols*}{2}

% TODO: invocations instead of spells

\subsection*{Rolling stats}

\paragraph{Virtues.} Creation of a witch starts the same as one for a knight and a spellweaver: roll \texttt{d6 + d12} four times; record not only the sum, but also individual values of each die, distribute them among \texttt{Virtues}, calculate \texttt{Focus} and \texttt{Stress}, determine \texttt{Injury slots}.

\paragraph{Guard.} Roll \texttt{d6} to determine the maximum \texttt{Guard(GD)} that witch has.

\paragraph{Choose a Passion. } Indulging a Passion restores \texttt{Spirit}, so choose it accordingly. 
That in large determines who what your witch actually draws their strength from.

\paragraph{Feats. } Witch knows no \texttt{Feats}. They had no access to martial art training or chose not to participate.

\newcolumn

\subsection*{Weapons \& Gear}

Without additional training, witches can't have more than one weapon dice per weapon they use(choose one if those are different), they also can't use armour plates.

\subsection*{Spellcasting}

Witch casts spells by invoking power of their patron. 
This is mutually beneficial agreement, patron expects for a witch to act on their behalf in the world of mortals.

Witch start with \texttt{2d4} favour and replenishes it by indulging patron's passions. 
At the start a witch knows 5 spells of their patron and all cantrips. 

\end{multicols*}
